\documentclass{article}
\usepackage{graphicx} % Required for inserting images

\title{aaaaa}
\author{Jordi Rodrígo Lozada Rivera}


\begin{document}
%RESUMEN
\begin{abstract}
En esta práctica se estudiaron dos fenómenos relacionados con el transporte de electrones en gases. En la primera parte se investigó el efecto Ramsauer--Townsend utilizando un tiratrón de xenón 2D21, registrando los voltajes de ánodo y de rejilla-blindaje como función del voltaje de cátodo para distintos voltajes de filamento. Las mediciones se realizaron tanto de manera manual como mediante un sistema automatizado, permitiendo comparar ambos métodos de adquisición de datos. En la segunda parte se empleó la técnica de Townsend pulsada para determinar la velocidad de deriva electrónica en una mezcla de 99\,\% Ar y 1\,\% N$_2$. A partir del tiempo de tránsito de un enjambre electrónico se calculó la velocidad de deriva para distintos valores del campo eléctrico reducido $E/N$. Los resultados obtenidos reproducen el comportamiento esperado para el transporte electrónico en gases y muestran la relación entre los procesos microscópicos de dispersión y las propiedades macroscópicas de transporte.
\end{abstract}

\vspace{0.3cm}

\textbf{Palabras clave:} efecto Ramsauer--Townsend, transporte electrónico, velocidad de deriva, técnica de Townsend pulsada, argón, xenón.

% RESULTADOS

\section{Resultados}

\subsection{Efecto Ramsauer--Townsend}

Las Figuras~5 y~6 presentan los voltajes de ánodo ($V_a$) y de rejilla-blindaje ($V_{rb}$) en función del voltaje de cátodo ($V_c$), obtenidos mediante adquisición automatizada y manual, respectivamente, para voltajes de filamento de $V_f=2$, $3$ y $4$ V.

En ambas modalidades de adquisición se observa que las curvas presentan la misma tendencia general para los tres valores de voltaje de filamento. Conforme aumenta $V_c$, tanto el voltaje de ánodo como el de rejilla experimentan variaciones continuas que reflejan los cambios en el número de electrones capaces de atravesar el gas contenido en el tiratrón.

Al incrementar el voltaje del filamento se produce una mayor emisión termoiónica desde el cátodo, por lo que la intensidad de la señal registrada aumenta ligeramente. Las mediciones obtenidas mediante el software presentan una menor dispersión respecto a las realizadas manualmente, lo que se atribuye a la eliminación del error asociado con la lectura de los instrumentos y con el ajuste manual del voltaje de cátodo.

En conjunto, ambos métodos reproducen el mismo comportamiento experimental, indicando que la automatización del proceso mejora la reproducibilidad sin modificar la respuesta física del sistema.

\subsection{Técnica de Townsend pulsada}

La Figura~7 muestra la velocidad de deriva electrónica $v_d$ en función del campo eléctrico reducido $E/N$ para las diferentes presiones estudiadas.

En todos los casos se observa que la velocidad de deriva depende del valor de $E/N$, lo que confirma que el transporte electrónico está gobernado principalmente por la energía adquirida por los electrones entre colisiones. Para valores bajos de $E/N$ predominan las colisiones con los átomos del gas y la velocidad de deriva permanece relativamente pequeña. Conforme el campo eléctrico reducido aumenta, los electrones adquieren mayor energía cinética entre colisiones y su velocidad media de desplazamiento también aumenta.

Las diferencias observadas entre las distintas presiones reflejan la variación en la frecuencia de colisiones electrón--átomo al modificarse la densidad del gas. Asimismo, el intervalo de valores de $E/N$ alcanzado para cada presión estuvo limitado por la aparición de descargas eléctricas dentro de la cámara experimental, condición a partir de la cual las mediciones dejaron de corresponder al régimen de corriente electrónica transitoria.

%% DISCUSIÓN Y CONCLUSIÓN


\section{Discusión y Conclusión}

Las mediciones realizadas con el tiratrón de xenón muestran que tanto el procedimiento manual como el automatizado producen resultados consistentes entre sí. Las diferencias observadas corresponden principalmente a la dispersión experimental inherente al método de adquisición manual, mientras que el sistema automatizado proporciona una mayor estabilidad y reproducibilidad de las mediciones.

El comportamiento registrado es compatible con el efecto Ramsauer--Townsend, fenómeno asociado con la disminución de la sección eficaz de dispersión electrónica para un intervalo específico de energías. Cuando los electrones presentan energías cercanas al mínimo de la sección eficaz, la probabilidad de colisión disminuye y una mayor fracción de ellos logra alcanzar el ánodo, modificando las señales eléctricas registradas en el dispositivo. Aunque el presente experimento no determina cuantitativamente la sección eficaz, sí permite observar sus consecuencias sobre el transporte electrónico en gases nobles.

En la segunda parte de la práctica, la técnica de Townsend pulsada permitió determinar la velocidad de deriva electrónica mediante la medición del tiempo de tránsito de un enjambre de electrones. La dependencia observada entre la velocidad de deriva y el campo eléctrico reducido concuerda con la teoría cinética del transporte de electrones, según la cual el cociente $E/N$ constituye el parámetro fundamental que determina la energía media adquirida por los electrones entre colisiones.

Las diferencias obtenidas para las distintas presiones son consecuencia del cambio en la densidad del gas y, por lo tanto, de la frecuencia de colisiones electrón--átomo. Asimismo, la aparición de descargas eléctricas estableció el límite superior de los valores de $E/N$ que pudieron explorarse experimentalmente.

En conclusión, ambas experiencias permitieron estudiar el transporte electrónico desde perspectivas complementarias. El experimento del efecto Ramsauer--Townsend puso de manifiesto la influencia de la sección eficaz de dispersión sobre el movimiento de electrones individuales, mientras que la técnica de Townsend pulsada permitió caracterizar el comportamiento colectivo de un enjambre electrónico mediante la determinación de su velocidad de deriva. Los resultados obtenidos fueron consistentes con la teoría revisada durante el curso y permitieron comprender la relación existente entre los procesos microscópicos de colisión y las propiedades macroscópicas del transporte de carga en gases.

\end{document}
